\section{The missing capabilities, not another redesign}
\label{sec:scope}

\subsection{The reviewed boundary}

This study uses Forge revision \code{674521027d96}, Leant revision
\code{6bf05ad78c46}, and standalone Djex revision \code{e8778f4ebd63}.
Leant's own Djex dependency is a different revision, \code{e237e866}; it must not
be silently replaced by the standalone head when reproducing Leant's reports.
The audit concentrates on Forge's merged structural-search and integration
sections, its provenance appendix, its explicit Leant/Djex review, and the
neighboring projects' current synthesis diagnostics and rank-N description
\cite{forge,forge-structural,forge-integration,forge-provenance,forge-neighbors,leant,djex-rank}.

The repository's design is already considerably broader than ``run a stronger
arithmetic tactic.'' In particular, induction generalization, demand-driven
lemma discovery, arithmetic witness synthesis, negative-evidence discipline,
and the separation of proposals from accepted proofs are already there.
Recommending them again would not answer the present request. The three
extensions below start precisely where those descriptions stop.

\begin{table}[ht]
\small
\begin{tabularx}{\textwidth}{@{}p{.23\textwidth}YY@{}}
\toprule
Area & Existing boundary & Addition developed here \\
\midrule
Higher-order synthesis & Type-directed candidates and higher-order carriers are available; whole compositions can still miss the budget. & Residual-specification-derived algebra contracts, local hole search, and a compatibility join before assembly. \\
Inductive equations & Polynomial recurrence synthesis and conserved quantities; general coupled ideal preservation is recognized as bilinear. & A descending linear-space algorithm that avoids joint bilinear synthesis, followed by target-directed ideal completion with polynomial multipliers. \\
Recursive proof search & Induction planning and acyclic tabling; a circular dependency is not a proof. & Typed cyclic obligation candidates compiled through explicit, checked well-founded descent on tagged states. \\
\bottomrule
\end{tabularx}
\caption{The contribution boundary. Existing material is used as infrastructure, not claimed as a new capability.}
\end{table}

The first column of additions is a search decomposition, the second is an
algebraic closure computation, and the third is a proof-construction strategy.
They are not proposed as new mathematical discoveries. Type-and-example-directed
synthesis, refinement-based synthesis, polynomial invariant analysis, and cyclic
proof systems are established research directions \cite{osera,polikarpova,karr,rodriguez,cycleq}.
The contribution is a concrete selection and adaptation for Forge, including
executable algorithms, certificate formats, failure cases, and integration gates.

\subsection{What the implementation actually establishes}

\status{Executed: finite-fragment Python synthesis; exact rational invariant-space
and ideal algorithms; ranking synthesis for supplied call graphs; an independent
certificate replayer; differential, mutation, and regression tests.
Not executed: a Lean tactic, source-to-polynomial reification, Lean kernel replay,
or a production change to Djex or Leant.}

The distinction is important because the strongest result here is not a claim
that a new \forge{} already beats \grind{}. There is no such end-to-end measurement.
The strongest demonstrated result is that small, precisely delimited workers can
produce structures unavailable to simpler ablations: a composition assembled from
locally constrained holes; a nonconserved but inductively closed family of
polynomials; and a decreasing interpretation of a cyclic call graph.

The source tree includes a proposed Lean file, \code{lean/Specimens.lean}, with
ordinary induction and well-founded-induction lemmas. It contains no \code{sorry},
but it was not compiled: no Lean executable was available in this execution
environment. Its presence is a starting point for porting, not a checking receipt.
This does not change the separate elaboration evidence already recorded by Forge
for its own core design files.

\subsection{A single operational principle}

All three extensions address the same practical failure: local search has found
useful pieces but has not discovered how they fit together. The remedy is to
make the required interface between pieces explicit.

For a fold, the interface is a base/constructor contract. For a family of
invariants, it is an action matrix or a matrix of polynomial multipliers. For
recursive obligations, it is a local proof constructor plus a ranking relation.
A successful search exports that interface, not just a Boolean success flag.

This yields a useful division of labor. Forge's existing obligation controller
owns the original theorem and invokes these workers. A worker may introduce
auxiliary obligations, but those obligations must have a mathematical reason:
a fold equation, an inductive closure condition, or a strict-descent requirement.
The controller does not need to guess arbitrary extra lemmas to make the worker
look productive.

The resulting development order should be: first exact proof lowering for the
algebraic certificates; then continuation-local synthesis with a checked fold
lemma; then the restricted cyclic compiler. The cyclic compiler depends on more
source-level structure than either algebraic lane, and should not be the first
component whose correctness is tested through a production tactic.
